\section{方法实现细节}

\subsection{TCForge 物理匹配合成平台}

本节详述主文~\S4.2 中合成数据平台 \eng{TCForge} 的实现，其目标是生成与真实 \eng{LWIR} 微扫描在物理特性上匹配的训练数据。

\subsubsection{场景几何生成}

\eng{TCForge} 通过 8 层随机原语叠加构造 \eng{IC} 布局的几何 mask：依次放置大块、中宽柱、L 形走线（25\% 粗 / 75\% 细）、框结构、细部与引脚阵列、圆形 via、对边 pin，最后做减通道处理。
所有层均由 seed 确定性控制。特征尺度由难度参数决定：easy 的 major/minor/line 为 80/60/45\,\(\mu\)m，medium 为 55/42/30\,\(\mu\)m，hard 为 35/28/22\,\(\mu\)m，stress 为 22/16/12\,\(\mu\)m。
场景整体旋转角为 \(\theta = 47.6° + U(-1.5°, +1.5°)\)，围绕实测 stage 角度抖动以匹配真实几何取向。

为消除阶梯锯齿对训练的污染，默认启用 4\(\times\)SSAA 抗锯齿：在 4 倍画布上绘制后旋转，再做 \(4 \times 4\) 块均值降采，得到 \(\text{coverage} \in [0, 1]\) 的连续值。
这样 \eng{HR} 目标不含阶梯锯齿，亚像素细线携带 coverage 缩放后的幅度。

\subsubsection{温度渲染与物理随机化}

温度场渲染公式为
\[
  T = T_{\mathrm{bg}} + \Delta T \cdot \text{coverage} + A \cdot \text{smooth\_noise}.
\]
关键物理参数的随机化分布经过仔细设计以覆盖实测范围：\(\Delta T \sim U(0.5, 5.0)\;^\circ\mathrm{C}\)（覆盖实测内/外轮廓对比度 1.94/2.49\,\(^\circ\)C）；
\eng{PSF} \(\sigma \sim U(0.15, 0.55)\) \eng{LR} px（覆盖采纳区间 [0.2, 0.5]，见 A.5.2），其中 30\% 使用椭圆 \eng{PSF}、10\% 使用 Airy 函数替代 Gaussian；
噪声遵循 lognormal 分布，均值等于实测噪声底 0.0724\,\(^\circ\)C；漂移幅度 \(\sim U(0, 0.3)\;^\circ\mathrm{C}\)；
亚像素位移使用实测 contour\_refined 的 248 帧 profile 重放并叠加 \(\sigma = 0.02\)\,px 的 jitter。

Forward 退化模式默认采用 physical\_block\_average：对 \eng{HR} 图先做 \eng{PSF} blur，再对各 shift 执行分块双线性均值降采（scale 2 或 4）。

\subsubsection{训练池}

当前构建了两个训练池。training\_pool\_2x\_aa 包含 1000 个 scene（难度分布 easy 200/medium 400/hard 300/stress 100），每 scene 248 帧，
\eng{LR} \(480 \times 640\) \(\rightarrow\) \eng{HR} \(960 \times 1280\)，用于 \(1\times\) 统计输入的各训练变体：\eng{Stats}、\eng{Stats+PixelShuffle}、\eng{Stats+HP-FC}、\eng{Stats+Full-FC}（run ids: v8.1a/v8.1b/v9b/v9d）。
training\_pool\_2x\_aa\_burst 在此基础上保存了完整的 \eng{LR} burst（约 152\,GB），用于 hybrid 输入的各变体：\eng{Hybrid}、\eng{Hybrid+Native-FC}、\eng{Hybrid+ResObs}（run ids: V9A/V9C/V10）。

为提高训练效率，还预计算了 \(K = 4\) 组 drizzle 变体：\(k = 0\) 使用全部 248 帧且无 shift 噪声（canonical，与推理一致）；
\(k = 1 \ldots 3\) 保留 60--100\% 帧（\(\ge 30\)）并叠加 \(N(0, 0.05\;\text{px})\) 的 shift 噪声。
训练时按 \([\text{seed}, \text{epoch}, \text{scene}, \text{0xBEEF}]\) 确定性抽样变体编号，将帧子集与 shift 噪声增广烘焙入预计算结果中，避免训练期间现场 drizzle 阻塞 dataloader。

\subsection{网络与损失}

\subsubsection{网络架构}

\eng{ThermalSRUNet} 采用编码-解码结构：编码路径为 3 层 MaxPool \(+\) ConvBlock（每块 \(2 \times 3 \times 3\) Conv \(+\) GroupNorm \(+\) SiLU \(+\) SE attention），
通道数按 \(c, 2c, 4c, 8c\) 递增（\(c = 64\)）；解码路径使用双线性上采样 \(+\) skip concat \(+\) ConvBlock。
\eng{HR} 输出头默认采用 bilinear 模式（双线性 \(\times\) scale \(\rightarrow\) \(3 \times 3\) Conv），
\eng{Stats+PixelShuffle}（v8.1b）消融变体另测试了 pixelshuffle 模式（ICNR 初始化 sub-pixel conv \(+\) PixelShuffle \(+\) HRResBlock），
后者产生条纹伪影且 proxy 全程更差（见 D.4.1），仅作为 head 归因对照保留。

在 hybrid 输入模式下，网络在 \(2\times\) 网格上以等效 scale\,=\,1 运行，因为输入已经是 \(2\times\) 分辨率。

\subsubsection{损失函数}

\eng{ContourSRLoss} 的总损失为
\[
  \mathcal{L} = w_m \cdot \text{MSE} + w_{hp} \cdot \text{HP} + w_e \cdot \text{Edge} + w_s \cdot (1 - \text{SSIM}) + w_{gv} \cdot \text{GV} + w_{lap} \cdot \text{Lap} + w_{fm} \cdot \text{FM} + w_r \cdot R.
\]
其中各项含义如下。MSE 为标准均方误差，可乘 gap\_weight。HP 为 highpass 结构损失 \(\|HP(\hat{x}) - HP(x)\|_1\)，预处理 \(HP(x) = x - G_{\sigma=5}(x)\)，乘以 structure/thin/gap 三种权重图。
Edge 为梯度幅值损失 \(\||\nabla\hat{x}| - |\nabla x|\|_1\) 加 0.25 倍半分辨率版本。SSIM 使用 \(11 \times 11\) Gaussian 窗。
GV（grad-vector）为 \(\|g_{x,p} - g_{x,t}\|_1 + \|g_{y,p} - g_{y,t}\|_1\)，保留梯度的完整向量信息以捕捉边缘膨胀和畸变——这是幅值型 Edge loss 看不到的方向偏差。
Lap 为 \(\text{ReLU}(|\Delta x| - |\Delta\hat{x}|)\) 的均值，只罚「比目标钝」。
FM 为 forward consistency \(\text{MSE}(\text{AvgPool}_s(G_{\sigma s}(\hat{x})),\; y_{\text{obs}})\)，可选 full-band 或 highpass 模式。
R 为残差惩罚 \(w_r \cdot \text{mean}(|\hat{x} - x_{\text{input}}|)\)（\eng{Hybrid+ResObs} / V10 使用）。

权重图机制使损失在结构关键区域有针对性地加强：structure\_boost 按梯度强度对 HP/GV 权重图加权；thin\_boost 在宽度 \(\le 3\) \eng{HR} px 的细结构上倍增；gap\_boost 在两侧有结构的窄缝背景上倍增。

conservative 权重壳（自 \eng{Stats} / v8.1a 起冻结，FC 与 hybrid 系列沿用）设定为 \(w_m = 0.3\)、\(w_{hp} = 0.8\)、structure\_boost \(=\) 2.0、\(w_{gv} = 0.15\)、thin\_boost \(=\) 3.0、gap\_boost \(=\) 2.0。
\eng{HotLoss}（v6）使用更激进的 hot 壳（\(w_{hp} = 1.0\)、structure\_boost \(=\) 4.0 等），叠加 Lap 0.1 和 full-band FM 0.1，保留为漂移放大参照变体。

\subsubsection{Hybrid 8 通道输入}

\eng{Hybrid}、\eng{Hybrid+Native-FC}、\eng{Hybrid+ResObs}（run ids: V9A/V9C/V10）使用 hybrid 8 通道输入，在 \(2\times\) 网格上融合 \(1\times\) 统计特征与 \(2\times\) drizzle 证据。
通道 0--4 为 \(1\times\) fused 统计（aligned mean/median/coverage/variance/highpass fused）经双线性上采到 \(2\times\)；
通道 5 为 drizzle mean @\(2\times\)，是观测域证据的主通道（也是 \eng{Hybrid+ResObs} 残差惩罚的基准）；
通道 6 为 drizzle coverage @\(2\times\)（空洞可识别：未观测 bin 以全局均值填充，coverage \(=\) 0）；通道 7 为 drizzle variance @\(2\times\)。

\eng{Hybrid+Native-FC}（V9C）的合法锚配套值得说明：hybrid 的通道 0 是上采样均值而非合法的 \(1\times\) 观测，
因此 forward 锚需另走数据管线携带原生 \(1\times\) aligned-mean patch（偶数 \(2\times\) origin 裁剪、增广同步）。

\subsection{训练配置对照}

七个变体的训练配置共享以下设定：scale\,=\,2、base\_channels\,=\,64、patch\_size\_hr\,=\,256、total\_steps\,=\,60000、lr\,=\,\(2 \times 10^{-4}\)、AMP、compile、
edge/ssim/coarse 权重 0.05/0.15/0.25、highpass\_sigma\,=\,5、real\_eval 使用 248 帧 contour\_refined \(+\) center-1/3 \(+\) zoom3。

各变体的差异字段构成一个 \(\{\text{input}\} \times \{\text{anchor}\}\) 消融矩阵。

\begin{table}[ht]
  \centering
  \caption{七个变体的差异字段对照}
  \label{tab:variant-config}
  \scriptsize
  \resizebox{\linewidth}{!}{%
  \begin{tabular}{lccccccc}
    \toprule
    字段 & HotLoss & Stats & Stats+PixelShuffle & Stats+HP-FC & Stats+Full-FC & Hybrid & Hybrid+Native-FC \\
    \midrule
    run id      & v6        & v8.1a    & v8.1b        & v9b      & v9d      & V9A             & V9C \\
    训练池      & 2x        & 2x\_aa   & 2x\_aa       & 2x\_aa   & 2x\_aa   & 2x\_aa\_burst   & 2x\_aa\_burst \\
    input / in\_ch & lr / 5 & lr / 5   & lr / 5       & lr / 5   & lr / 5   & hybrid / 8      & hybrid / 8 \\
    HR 头       & bilinear  & bilinear & pixelshuffle & bilinear & bilinear & bilinear        & bilinear \\
    loss 壳     & hot+Lap   & cons.    & cons.        & cons.    & cons.    & cons.           & cons. \\
    forward 锚  & full 0.1  & 无       & 无           & hp 0.1   & full 0.1 & 无              & legal-1x hp 0.1 \\
    batch\_size & 128       & 128      & 128          & 128      & 128      & 64              & 64 \\
    \bottomrule
  \end{tabular}
  }
\end{table}

消融矩阵的直读方式：\(\{1\times, \text{hybrid}\} \times \{\text{none}, \text{HP-FC}, \text{Full-FC}, \text{Native-FC}\}\) \(=\) \eng{Stats} / \eng{Stats+HP-FC} / \eng{Stats+Full-FC} / \eng{Hybrid} / \eng{Hybrid+Native-FC}（\eng{HotLoss} 和 \eng{Stats+PixelShuffle} 分别为 loss 温度和 head 的额外归因变体）。

\textbf{可比性说明：} \eng{Hybrid}（V9A）前 35K 步使用 batch\_size\,=\,128，之后与 \eng{Hybrid+Native-FC}（V9C）全程 batch\_size\,=\,64——\eng{Hybrid} 与 \eng{Hybrid+Native-FC} 两个变体与 \(1\times\) 各变体的训练动力学不完全同条件。
\eng{Hybrid}（V9A）30K 保真悬崖与 batch size 切换时间重合，混杂未排除。

\subsection{经典方法}

\subsubsection{Drizzle}

Drizzle 使用 STScI drizzle 库实现，pixfrac 默认 0.7（sweep \{1.0, 0.8, 0.7, 0.6, 0.5\}），kernel 为 square，输出网格 \((H \times 2, W \times 2)\)。
coverage \(< 1.0\) 的像素置 NaN。\eng{TCForge} 训练侧的 scatter 使用 bilinear kernel，两种 kernel 的差异在 D.7 的 fine-window 指标中以 drizzle 输入通道参照点的形式体现。

\subsubsection{各向异性 TGV}

\eng{TGV} 重建采用外层 FISTA 框架 \(x_{k+1} = \text{TGV\_prox}(z_k - \eta \nabla D(z_k))\)，针对本数据集的 raster 各向异性（见 B.3.1）做了两项关键设计。

\textbf{第一，各向异性对偶投影。} \eng{TGV} 内层（Chambolle-Pock 路径）的一阶对偶球改为椭圆 \(\{(a,b) : (a/r_a)^2 + (b/r_b)^2 \le 1\}\)，
其中 \(r_a = \alpha_1 \cdot r_y\)、\(r_b = \alpha_1\)（\(r_y = 1.5\)）；二阶对称张量同构造。这使得 Y 方向正则强度弱于 X 方向，补偿 Y 方向数据约束的不足。

\textbf{第二，coverage 加权数据梯度。} 每个 \eng{HR} 像素的数据保真梯度除以预计算的 coverage（由 bilinear splat \(+\) \eng{PSF} adjoint 得到），而非统一除以帧数 \(N\)。
这修正了 bilinear scatter 把权重集中在固定 \eng{HR} 行上所导致的水平条纹。

参数网格为 \(\lambda_{\text{tv}} \in \{3 \times 10^{-4}, 10^{-3}, 3 \times 10^{-3}\} \times \sigma_{\text{PSF}} \in \{0.18, 0.50\}\)，\(\alpha_{\text{ratio}} = 2.0\)（\(\alpha_1 = \lambda\), \(\alpha_0 = \lambda/2\)），外层 100 iter、内层 80 iter。
最优参数下 artifact 从 3.870 降至 0.695（\(-82\%\)），raw-control corr 从 0.902 升至 0.916，耗时 30.8\,min CPU。

\subsubsection{MAP-TV 去卷积基准}

\eng{MAP-TV} 基准用于验收门控（C.5.0 采纳规则）：学习方法必须在 \eng{FRC} 频带一致性与 zigzag profile 指标上同时不输 \eng{MAP-TV} 才可采纳。
其 forward 模型为 shift \(\rightarrow\) Gaussian \eng{PSF}（\(\sigma\) \eng{LR} px）\(\rightarrow\) avg\_pool box，使用 GPU FISTA \(+\) smoothed-TV 梯度优化，150 iter。

参数网格为 \(\sigma \in \{0.2, 0.3, 0.4, 0.5\} \times \lambda \in \{3 \times 10^{-4}, 10^{-3}, 3 \times 10^{-3}\}\)。
选择规则为每个 \(\sigma\) 内按 split-half proxy（split\_half\_nrmse \(+\) 0.05\(\times\)artifact \(+\) 0.08\(\times\)std\_excess）取最小，再跨 \(\sigma\) 选全局最优，结果为 \(\sigma = 0.2\)、\(\lambda = 10^{-3}\)，耗时 4563\,s GPU。

\subsection{Checkpoint 选择协议}

训练期间的 checkpoint 选择对学习变体的最终输出质量至关重要。本节详述主文~\S4.5 和~\S5.6 的选择规则与采纳判定。

\subsubsection{采纳规则（Adoption rule）}

一个学习方法被采纳，当且仅当在其选定 checkpoint 上同时满足：
（a）\eng{FRC} 频带一致性 \(\ge\) \eng{MAP-TV} 基准；（b）zigzag FWHM/dip 不差于基准；（c）proxy 轨迹无 post-selection 漂移悬崖；（d）双域视觉 panel 均通过。
否则经典基准仍是交付物。正是这条保守规则，让诚实的负结果（PixelShuffle 头、\(4\times\) 网络、loss 侧锚定）与阳性结果得以并列报告——拒绝一个方法和采纳一个方法用的是同一把尺子。

\subsubsection{选择流程}

对每个变体的所有 eval checkpoint，首先归一化两个 proxy：\(a_{\text{norm}} = (\text{artifact} - \min) / (\max - \min)\)，\(c_{\text{norm}} = (\max_{\text{corr}} - \text{corr}) / (\max_{\text{corr}} - \min_{\text{corr}})\)（corr 反向）。
计算到理想点 \((0, 0)\) 的距离 \(d = \sqrt{a_{\text{norm}}^2 + c_{\text{norm}}^2}\)，按 \(d\) 升序遍历。
为避免选中相邻 checkpoint 的冗余候选，采用 5K 窗口去重：每个 checkpoint 归属 window \(=\) step // 5000，同窗只保留最先入选者，最多保留 3 个候选。
60K 步若未入选则强制追加为漂移参照。Rank-1 为 canonical 候选。最终每个候选需通过温度域三联 panel 的视觉 gate，人工过门后定稿。

\subsubsection{已执行选点}

\eng{EP11} 四个变体的选点结果如下：\eng{HotLoss}（v6）canonical 8K（artifact/corr \(=\) 0.330/0.774），\eng{Stats}（v8.1a）canonical 15K（0.392/0.758），
\eng{Stats+PixelShuffle}（v8.1b）canonical 5K（0.370/0.739），\eng{Stats+HP-FC}（v9b）canonical 11K（0.339/0.777）。\eng{TGV} Pareto 参考点为 (0.695, 0.916)。
各变体 60K 端点的 artifact 均显著恶化——按端点上报会让每个变体交出最差 checkpoint，这正是选择协议存在的意义。

统一 harness 后的补充选点为：\eng{Hybrid}（V9A）canonical 10K（harness artifact/corr \(=\) 1.762/0.719），\eng{Hybrid+Native-FC}（V9C）canonical 5K（1.669/0.718），
\eng{Stats+Full-FC}（V9D）canonical 7K（1.726/0.771），\eng{Hybrid+ResObs}（V10）采用高-\(\lambda\) sweep 的工作点 \(\lambda=1.2\)@15K（2.726/0.711；fine-window 为 hp\_corr\_input 0.922、sharp\_p95 0.987、lattice 0.0141）。
\eng{Hybrid}（V9A）60K 不作为交付 checkpoint，仅作为 F5 late-drift visual control。
需强调 hybrid 变体与 \(1\times\) 变体的 proxy 不可跨列横比（见 A.3.2）；最终 T1/T2 横评只使用 \texttt{output/ep11\_unified\_harness/} 的同一 harness scale。

\subsection{结构级指标与视觉门控}

本节详述主文实验中使用的结构级评估指标与视觉门控流程。

\subsubsection{Zigzag 走线 profile 指标}

以客户关心的中心走线为测量对象。在固定横截面上，对每条 profile 测表观 FWHM（full width at half maximum）与谷深（dip），报告 median FWHM/dip 及 per-profile 分布。
当 per-profile 结果混杂时如实报告为「混杂」，不取巧。需要注意，\(2\times\) 网格下量化步长为 5\,\(\mu\)m，FWHM 可能饱和（主文~\S5.1 多个方法同为 40\,\(\mu\)m），区分度低时不应过度解读。

\subsubsection{双域视觉 panel}

双域 panel 由两张并列图组成。第一张为高通结构图（边缘证据），使用 \(HP(x) = x - G_{\sigma=5}(x)\)，红/蓝分别表示相对局部背景的正/负有符号响应，白色表示无变化。
第二张为普通温度图（inferno colormap，1--99 百分位裁剪），专门用于捕捉「只有边缘变亮但温度场整体失真」这类被高通图美化掉的失效模式。
两图并列能让审查者同时看到高频结构证据和低频温度合理性。

\subsubsection{对齐门控一致性}

重建不得劣化被门控锚点上的 hold-out 轮廓一致性。该角色由 \eng{EP04} 对齐质量评估承担：
若某方法在门控锚点上的 split-half 轮廓一致性低于对照，该方法在该区域不予采纳。
